% Active characters, math-active characters and catcodes: the same tokens read another way.
\documentclass{article}
\usepackage{amsmath,amssymb}
\usepackage{tikz-cd}
\pdfpagewidth=5in \pdfpageheight=5in
%%%% case babel-french
%%%% pre
\usepackage[T1]{fontenc}
\usepackage[french]{babel}
%%%% body
\begin{tikzcd} A : B \arrow[r, "f"] & C ; D \arrow[d, "g : h"] \\ & E ! F ? \end{tikzcd}
%%%% case babel-french-library
%%%% pre
\usepackage[T1]{fontenc}
\usepackage[french]{babel}
\usetikzlibrary{babel}
%%%% body
\begin{tikzcd} A : B \arrow[r, "f"] & C ; D \arrow[d, "g : h"] \\ & E ! F ? \end{tikzcd}
%%%% case babel-french-plain
%%%% pre
\usepackage[T1]{fontenc}
\usepackage[french]{babel}
%%%% body
\begin{tikzcd} A \arrow[r, "f"] & B \end{tikzcd}
%%%% case babel-french-text
%%%% pre
\usepackage[T1]{fontenc}
\usepackage[french]{babel}
%%%% body
\begin{tikzcd} \text{a : b ; c ! d ?} \arrow[r] & \mbox{x : y} \end{tikzcd}
%%%% case babel-french-picture-after
%%%% pre
\usepackage[T1]{fontenc}
\usepackage[french]{babel}
%%%% body
\begin{tikzcd} A \arrow[r] & B \end{tikzcd}
Du texte : et ceci ; et cela ! et quoi ?
\begin{tikzcd} C \arrow[r] & D \end{tikzcd}
%%%% case babel-french-off
%%%% pre
\usepackage[T1]{fontenc}
\usepackage[french]{babel}
%%%% body
\shorthandoff{:;!?}
\begin{tikzcd} A : B \arrow[r, "f"] & C ; D \end{tikzcd}
%%%% case babel-german
%%%% pre
\usepackage[T1]{fontenc}
\usepackage[german]{babel}
%%%% body
\begin{tikzcd} A \arrow[r, "f"] & B \end{tikzcd}
%%%% case babel-ngerman
%%%% pre
\usepackage[T1]{fontenc}
\usepackage[ngerman]{babel}
%%%% body
\begin{tikzcd} A \arrow[r] & B \end{tikzcd}
%%%% case babel-german-library
%%%% pre
\usepackage[T1]{fontenc}
\usepackage[ngerman]{babel}
\usetikzlibrary{babel}
%%%% body
\begin{tikzcd} A \arrow[r, "f"] & B \end{tikzcd}
%%%% case babel-german-shorthand-cell
%%%% pre
\usepackage[T1]{fontenc}
\usepackage[ngerman]{babel}
\usetikzlibrary{babel}
%%%% body
\begin{tikzcd} \text{"a "o} \arrow[r] & B \end{tikzcd}
%%%% case babel-german-off
%%%% pre
\usepackage[T1]{fontenc}
\usepackage[ngerman]{babel}
%%%% body
\shorthandoff{"}
\begin{tikzcd} A \arrow[r, "f"] & B \end{tikzcd}
\shorthandon{"}
%%%% case babel-spanish
%%%% pre
\usepackage[T1]{fontenc}
\usepackage[spanish]{babel}
%%%% body
\begin{tikzcd} A \arrow[r, "f"] & B < C > D \end{tikzcd}
%%%% case babel-spanish-library
%%%% pre
\usepackage[T1]{fontenc}
\usepackage[spanish]{babel}
\usetikzlibrary{babel}
%%%% body
\begin{tikzcd} A \arrow[r, "f"] & B < C > D ~ E \end{tikzcd}
%%%% case babel-spanish-math-dot
%%%% pre
\usepackage[T1]{fontenc}
\usepackage[spanish,es-nodecimaldot]{babel}
%%%% body
\begin{tikzcd} 1.5 \arrow[r] & 2.5 \end{tikzcd}
%%%% case babel-turkish
%%%% pre
\usepackage[T1]{fontenc}
\usepackage[turkish]{babel}
%%%% body
\begin{tikzcd} A = B \arrow[r, "x = y"] & C \end{tikzcd}
%%%% case babel-turkish-library
%%%% pre
\usepackage[T1]{fontenc}
\usepackage[turkish]{babel}
\usetikzlibrary{babel}
%%%% body
\begin{tikzcd}[row sep=1em] A = B \arrow[r, "x = y"] & C \end{tikzcd}
%%%% case active-letter-arrow
\catcode`\Q=13 \def Q{\arrow[r]}
\begin{tikzcd} A Q & B \end{tikzcd}
%%%% case active-semicolon-arrow
\catcode`\;=13 \def;{\arrow[r]}
\begin{tikzcd} A ; & B \end{tikzcd}
%%%% case active-colon-text
\catcode`\:=13 \def:{\mathrm{z}}
\begin{tikzcd} A : \arrow[r] & B \end{tikzcd}
%%%% case active-bar-text
\catcode`\|=13 \def|{\mathrm{z}}
\begin{tikzcd} A | \arrow[r] & B \end{tikzcd}
%%%% case active-bar-first
\catcode`\|=13 \def|{\mathrm{z}}
\begin{tikzcd} | A \arrow[r] & B \end{tikzcd}
%%%% case active-bar-options
\catcode`\|=13 \def|{\mathrm{z}}
\begin{tikzcd} |[draw]| A \arrow[r] & B \end{tikzcd}
%%%% case active-bar-eats
\catcode`\|=13 \def|#1{\mathrm{#1}}
\begin{tikzcd} |x \arrow[r] & B \end{tikzcd}
%%%% case active-comma-in-keys
\catcode`\,=13 \def,{\mathrm{z}}
\begin{tikzcd} A \arrow[r, "f"] & B \end{tikzcd}
%%%% case active-comma-in-cell
\catcode`\,=13 \def,{\mathrm{z}}
\begin{tikzcd} A , \arrow[r] & B \end{tikzcd}
%%%% case active-eq-in-keys
\catcode`\==13 \def={\mathrm{z}}
\begin{tikzcd} A \arrow[r, shift left=1ex] & B \end{tikzcd}
%%%% case active-bracket-open
\catcode`\[=13 \def[{\mathrm{z}}
\begin{tikzcd} A & B \end{tikzcd}
%%%% case active-bracket-keys
\catcode`\]=13 \def]{\mathrm{z}}
\begin{tikzcd} A \arrow[r] & B \end{tikzcd}
%%%% case active-quote-in-label
\catcode`\"=13 \def"{\mathrm{z}}
\begin{tikzcd} A \arrow[r, "f"] & B \end{tikzcd}
%%%% case active-quote-brace
\catcode`\'=13 \def'{\mathrm{z}}
\begin{tikzcd} A \arrow[r, "f"'] & B \end{tikzcd}
%%%% case active-tilde
\begin{tikzcd} A~B \arrow[r] & C \end{tikzcd}
%%%% case active-tilde-redefined
\def~{\arrow[r]}
\begin{tikzcd} A~ & C \end{tikzcd}
%%%% case active-ampersand
\catcode`\&=13 \def&{\mathrm{z}}
\begin{tikzcd} A & B \end{tikzcd}
%%%% case other-ampersand
\catcode`\&=12
\begin{tikzcd} A & B \end{tikzcd}
%%%% case letter-ampersand
\catcode`\&=11
\begin{tikzcd} A & B \end{tikzcd}
%%%% case ampersand-another-tab
\catcode`\&=12 \catcode`\Q=4
\begin{tikzcd} A Q B \end{tikzcd}
%%%% case backslash-another
\catcode`\|=0
|begin{tikzcd} A |arrow[r] & B |end{tikzcd}
%%%% case backslash-another-row
\catcode`\|=0
|begin{tikzcd} A |arrow[d] |\ B |end{tikzcd}
%%%% case brace-another
\catcode`\<=1 \catcode`\>=2
\begin{tikzcd} \mathrm<x> \arrow[r] & B \end{tikzcd}
%%%% case brace-another-key
\catcode`\<=1 \catcode`\>=2
\begin{tikzcd} A \arrow[r, "f"<above>] & B \end{tikzcd}
%%%% case math-shift-another
\catcode`\Q=3
\begin{tikzcd} \text{Qx^2Q} \arrow[r] & B \end{tikzcd}
%%%% case sup-other
\catcode`\^=12
\begin{tikzcd} A^B \arrow[r] & B \end{tikzcd}
%%%% case sub-other
\catcode`\_=12
\begin{tikzcd} A_B \arrow[r] & B \end{tikzcd}
%%%% case sup-letter-key
\catcode`\^=11
\begin{tikzcd} A \arrow[r, "f"^] & B \end{tikzcd}
%%%% case sup-active
\catcode`\^=13 \def^{\mathrm{z}}
\begin{tikzcd} A^B \arrow[r] & B \end{tikzcd}
%%%% case sup-mark-another
\catcode`\Q=7
\begin{tikzcd} AQB \arrow[r] & B \end{tikzcd}
%%%% case sub-mark-another
\catcode`\Q=8
\begin{tikzcd} AQB \arrow[r] & B \end{tikzcd}
%%%% case space-other
\catcode`\ =12
\begin{tikzcd} A \arrow[r] & B \end{tikzcd}
\catcode`\ =10
%%%% case space-active
\catcode`\ =13 \def {\mathrm{z}}%
\begin{tikzcd} A \arrow[r] & B \end{tikzcd}%
\catcode`\ =10
%%%% case space-letter-math-active
\catcode`\ =12 \mathcode`\ ="8000%
\begin{tikzcd} A \arrow[r] & B \end{tikzcd}%
\catcode`\ =10
%%%% case obeyspaces
\obeyspaces
\begin{tikzcd} A \arrow[r] & B \end{tikzcd}
%%%% case obeylines
\obeylines
\begin{tikzcd} A \arrow[r] & B \end{tikzcd}
%%%% case obeylines-multi
{\obeylines
\begin{tikzcd}
A \arrow[r] & B
\end{tikzcd}
}
%%%% case endlinechar-none
\endlinechar=-1
\begin{tikzcd}
A \arrow[r] & B
\end{tikzcd}
\endlinechar=13
%%%% case endlinechar-letter
\endlinechar=`\Q
\begin{tikzcd}
A \arrow[r] & B
\end{tikzcd}
\endlinechar=13
%%%% case carriage-other
\catcode13=12
\begin{tikzcd}
A \arrow[r] & B
\end{tikzcd}\catcode13=5
%%%% case comment-another
\catcode`\Q=14
\begin{tikzcd} A \arrow[r] & B Q & C
\end{tikzcd}
%%%% case percent-other
\catcode`\%=12
\begin{tikzcd} A \arrow[r] & 50% \end{tikzcd}
%%%% case percent-active
\catcode`\%=13 \def%{\mathrm{z}}
\begin{tikzcd} A \arrow[r] & 50% \end{tikzcd}
\catcode`\%=14
%%%% case hash-other
\catcode`\#=12
\begin{tikzcd} A # \arrow[r] & B \end{tikzcd}
%%%% case hash-active
\catcode`\#=13 \def#{\mathrm{z}}
\begin{tikzcd} A # \arrow[r] & B \end{tikzcd}
%%%% case dollar-other
\catcode`\$=12
\begin{tikzcd} A $ \arrow[r] & B \end{tikzcd}
\catcode`\$=3
%%%% case ignored-char
\catcode`\Q=9
\begin{tikzcd} AQB \arrow[r] & B \end{tikzcd}
%%%% case invalid-char
\catcode`\Q=15
\begin{tikzcd} AQB \arrow[r] & B \end{tikzcd}
%%%% case letter-other-r
\catcode`\r=12
\begin{tikzcd} A \arrow[r] & B \end{tikzcd}
%%%% case letter-other-r-separated
{\catcode`\r=12 \gdef\rother{r}}
\begin{tikzcd} A \arrow[\rother] & B \end{tikzcd}
%%%% case letter-other-key
{\catcode`\r=12 \gdef\x{\begin{tikzcd} A \arrow[r] \\ B \end{tikzcd}}}
\x
%%%% case letter-other-d
\catcode`\d=12
\begin{tikzcd} A \arrow[d] \\ B \end{tikzcd}
%%%% case letter-digit
\catcode`\1=11
\begin{tikzcd}[row sep=1cm] A \arrow[d] \\ B \end{tikzcd}
%%%% case letter-at-other
\makeatother
\begin{tikzcd} A@B \arrow[r] & B \end{tikzcd}
%%%% case letter-at-letter
\makeatletter
\begin{tikzcd} A@B \arrow[r] & B \end{tikzcd}
%%%% case letter-at-macro
\makeatletter
\begin{tikzcd} A\@empty B \arrow[r] & B \end{tikzcd}
%%%% case hat-notation-letter
\begin{tikzcd} A^^41 \arrow[r] & B \end{tikzcd}
%%%% case hat-notation-amp
\begin{tikzcd} A ^^26 B \end{tikzcd}
%%%% case hat-notation-cs
\begin{tikzcd} A \arro^^77[r] & B \end{tikzcd}
%%%% case hat-notation-newline
\begin{tikzcd} A ^^M \arrow[r] & B \end{tikzcd}
%%%% case high-char
\begin{tikzcd} \text{caf^^e9} \arrow[r] & B \end{tikzcd}
%%%% case inputenc-latin1
%%%% pre
\usepackage[latin1]{inputenc}
%%%% body
\begin{tikzcd} \text{é} \arrow[r] & B \end{tikzcd}
%%%% case inputenc-utf8-text
\begin{tikzcd} \text{é} \arrow[r, "\text{ü}"] & \text{ñ} \end{tikzcd}
%%%% case inputenc-utf8-math
\begin{tikzcd} é \arrow[r] & B \end{tikzcd}
%%%% case math-active-letter-arrow
\mathcode`\Q="8000 {\catcode`\Q=13 \gdef Q{\arrow[r]}}
\begin{tikzcd} A Q & B \end{tikzcd}
%%%% case math-active-letter-through-macro
\mathcode`\Q="8000 {\catcode`\Q=13 \gdef Q{\arrow[r]}}
\newcommand\h{Q}
\begin{tikzcd} A \h & B \end{tikzcd}
%%%% case math-active-letter-amp-through-macro
\mathcode`\Q="8000 {\catcode`\Q=13 \gdef Q{\\}}
\newcommand\h{Q}
\begin{tikzcd} A \h B \end{tikzcd}
%%%% case math-active-letter-in-body-only
\mathcode`\Q="8000 {\catcode`\Q=13 \gdef Q{\mathrm{z}}}
\newcommand\h{Q}
\begin{tikzcd} A \h \arrow[r] & B \end{tikzcd}
%%%% case math-active-prime
\begin{tikzcd} A' \arrow[r] & B'' \end{tikzcd}
%%%% case math-active-prime-redefined
\begingroup\catcode`\'=13 \gdef'{\arrow[r]}\endgroup
\begin{tikzcd} A' & B \end{tikzcd}
%%%% case math-active-prime-redefined-macro
\begingroup\catcode`\'=13 \gdef'{\arrow[r]}\endgroup
\newcommand\h{A'}
\begin{tikzcd} \h & B \end{tikzcd}
%%%% case math-active-prime-label
\begin{tikzcd} A \arrow[r, "f'"] & B \end{tikzcd}
%%%% case math-active-prime-swap
\begin{tikzcd} A \arrow[r, "f'"'] & B \end{tikzcd}
%%%% case math-active-char
\mathcode`\Q="8000 {\catcode`\Q=13 \gdef Q{\arrow[r]}}
\begin{tikzcd} A \char`\Q & B \end{tikzcd}
%%%% case math-active-chardef
\mathcode`\Q="8000 {\catcode`\Q=13 \gdef Q{\arrow[r]}}
\chardef\q=`\Q
\begin{tikzcd} A \q & B \end{tikzcd}
%%%% case math-active-chardef-macro
\mathcode`\Q="8000 {\catcode`\Q=13 \gdef Q{\arrow[r]}}
\chardef\q=`\Q
\newcommand\h{\q}
\begin{tikzcd} A \h & B \end{tikzcd}
%%%% case math-active-let-char
\mathcode`\Q="8000 {\catcode`\Q=13 \gdef Q{\arrow[r]}}
\let\q=Q
\begin{tikzcd} A \q & B \end{tikzcd}
%%%% case math-active-digit
\mathcode`\1="8000 {\catcode`\1=13 \gdef 1{\arrow[r]}}
\begin{tikzcd} A 1 & B \end{tikzcd}
%%%% case math-active-digit-key
\mathcode`\1="8000 {\catcode`\1=13 \gdef 1{\mathrm{z}}}
\begin{tikzcd} A \arrow[r, shift left=1pt] & B \end{tikzcd}
%%%% case math-active-digit-label
\mathcode`\1="8000 {\catcode`\1=13 \gdef 1{\mathrm{z}}}
\begin{tikzcd} A \arrow[r, "1"] & B 1 \end{tikzcd}
%%%% case math-active-minus
\mathcode`\-="8000 {\catcode`\-=13 \gdef-{\mathrm{z}}}
\begin{tikzcd} A - B \arrow[r] & C \end{tikzcd}
%%%% case math-active-in-text
\mathcode`\Q="8000 {\catcode`\Q=13 \gdef Q{\arrow[r]}}
\begin{tikzcd} \text{Q} \arrow[r] & B \end{tikzcd}
%%%% case math-active-space
\mathcode`\ ="8000 {\catcode`\ =13 \gdef {\arrow[r]}}
\begin{tikzcd} A B \arrow[d] & C \\ D \end{tikzcd}
%%%% case math-active-undefined
\mathcode`\Q="8000 {\catcode`\Q=13 \global\let Q\undefined}
\begin{tikzcd} A Q \arrow[r] & B \end{tikzcd}
%%%% case declare-math-symbol
\DeclareMathSymbol{\foo}{\mathrel}{symbols}{"20}
\begin{tikzcd} A \foo B \arrow[r] & C \end{tikzcd}
%%%% case declare-math-symbol-redefine-alpha
\DeclareMathSymbol{\alpha}{\mathalpha}{letters}{"0C}
\begin{tikzcd} \alpha \arrow[r] & C \end{tikzcd}
%%%% case declare-math-symbol-letter
\DeclareMathSymbol{Q}{\mathrel}{symbols}{"20}
\begin{tikzcd} A Q B \arrow[r] & C \end{tikzcd}
%%%% case declare-math-active-symbol
\DeclareMathSymbol{\foo}{\mathalpha}{letters}{"0C}
\mathcode`\Q=\foo
\begin{tikzcd} A Q B \arrow[r] & C \end{tikzcd}
%%%% case declare-robust-symbol
\newcommand\foo{\mathrel{\arrow[r]}}
\begin{tikzcd} A \foo B & C \end{tikzcd}
%%%% case mathcode-digit-var
\mathcode`\1="7031
\begin{tikzcd} 1 \arrow[r] & 2 \end{tikzcd}
%%%% case delcode
\delcode`\/="02F30E
\begin{tikzcd} \left/ A \right/ \arrow[r] & 2 \end{tikzcd}
%%%% case lccode-active
\begingroup\lccode`\~=`\Q \lowercase{\endgroup\def~{\arrow[r]}}
\catcode`\Q=13
\begin{tikzcd} A Q & B \end{tikzcd}
%%%% case lowercase-active-in-macro
\begingroup\lccode`\~=`\&\lowercase{\endgroup\def\h{~}}
\begin{tikzcd} A \h B \end{tikzcd}
%%%% case uppercase-direct
\begin{tikzcd} \uppercase{a} \arrow[r] & B \end{tikzcd}
%%%% case uppercase-arrow
\begin{tikzcd} \uppercase{\arrow[r]} A & B \end{tikzcd}
%%%% case uppercase-key-chars
\uccode`\[=`\[
\begin{tikzcd} A \uppercase{\ar}[r] & B \end{tikzcd}
%%%% case lccode-amp
\lccode`\Q=`\&
\newcommand\h{\lowercase{Q}}
\begin{tikzcd} A \h B \end{tikzcd}
%%%% case catcode-in-cell
\begin{tikzcd} A \catcode`\Q=12 \arrow[r] & B \end{tikzcd}
%%%% case catcode-in-macro
\newcommand\h{\catcode`\Q=13 }
\begin{tikzcd} A \h \arrow[r] & Q \end{tikzcd}
%%%% case catcode-let
\let\h\catcode
\begin{tikzcd} A \h`\Q=13 \arrow[r] & Q \end{tikzcd}
%%%% case makeatletter-in-cell
\begin{tikzcd} A \makeatletter \arrow[r] & B@C \end{tikzcd}
%%%% case scantokens-in-cell
\begin{tikzcd} A \scantokens{x} \arrow[r] & B \end{tikzcd}
%%%% case scantokens-in-macro
\newcommand\h{\scantokens{\arrow[r]}}
\begin{tikzcd} A \h & B \end{tikzcd}
%%%% case verb-in-cell
\begin{tikzcd} A \verb|x| \arrow[r] & B \end{tikzcd}
%%%% case verb-hides-amp
\begin{tikzcd} A \verb|&| \arrow[r] & B \end{tikzcd}
%%%% case detokenize-amp
\begin{tikzcd} A \detokenize{&} \arrow[r] & B \end{tikzcd}
%%%% case detokenize-arrow
\begin{tikzcd} A \detokenize{\arrow[r]} & B \end{tikzcd}
%%%% case string-arrow
\begin{tikzcd} A \string\arrow[r] & B \end{tikzcd}
%%%% case string-amp
\begin{tikzcd} A \string& B \end{tikzcd}
%%%% case string-row
\begin{tikzcd} A \string\\ B \end{tikzcd}
%%%% case meaning-arrow
\begin{tikzcd} A \meaning\arrow[r] & B \end{tikzcd}
%%%% case meaning-macro
\newcommand\h{x}
\begin{tikzcd} \meaning\h \arrow[r] & B \end{tikzcd}
%%%% case number-count
\begin{tikzcd} \number\count255 \arrow[r] & B \end{tikzcd}
%%%% case inputlineno
\begin{tikzcd} \the\inputlineno \arrow[r] & \the\inputlineno \\
\the\inputlineno & \the\inputlineno
\end{tikzcd}
%%%% case inputlineno-multiline-cell
\begin{tikzcd} \the\inputlineno
\the\inputlineno
\the\inputlineno \arrow[r] & B \end{tikzcd}
%%%% case currentgrouplevel
\begin{tikzcd} \the\currentgrouplevel \arrow[r] & B \end{tikzcd}
%%%% case currentgrouptype
\begin{tikzcd} \the\currentgrouptype \arrow[r] & B \end{tikzcd}
%%%% case currentiflevel
\begin{tikzcd} \the\currentiflevel \arrow[r] & B \end{tikzcd}
%%%% case currentiftype
\begin{tikzcd} \the\currentiftype \arrow[r] & B \end{tikzcd}
%%%% case currentifbranch
\begin{tikzcd} \the\currentifbranch \arrow[r] & B \end{tikzcd}
%%%% case lastnodetype
\begin{tikzcd} \the\lastnodetype \arrow[r] & B \end{tikzcd}
%%%% case lastpenalty
\begin{tikzcd} \penalty7 \the\lastpenalty \arrow[r] & \the\lastkern \the\lastskip \end{tikzcd}
%%%% case badness
\begin{tikzcd} \the\badness \arrow[r] & B \end{tikzcd}
%%%% case interactionmode
\begin{tikzcd} \the\interactionmode \arrow[r] & B \end{tikzcd}
%%%% case hsize
\begin{tikzcd} \the\hsize \arrow[r] & \the\linewidth \end{tikzcd}
%%%% case parindent
\begin{tikzcd} \the\parindent \arrow[r] & \the\parskip \end{tikzcd}
%%%% case tabskip
\begin{tikzcd} \the\tabskip \arrow[r] & B \end{tikzcd}
%%%% case tabskip-document
\tabskip=3pt
\begin{tikzcd} \the\tabskip \arrow[r] & B \end{tikzcd}
%%%% case baselineskip
\begin{tikzcd} \the\baselineskip \arrow[r] & B \end{tikzcd}
%%%% case baselineskip-parbox
\begin{tikzcd} \parbox{1cm}{\the\baselineskip} \arrow[r] & B \end{tikzcd}
%%%% case baselinestretch
\renewcommand\baselinestretch{2}\selectfont
\begin{tikzcd} \parbox{1cm}{a b c d e f g h i j k l m n} \arrow[r] & B \end{tikzcd}
%%%% case lineskip
\lineskip=5pt \lineskiplimit=20pt
\begin{tikzcd} \the\lineskip \arrow[r] & \the\lineskiplimit \end{tikzcd}
%%%% case spaceskip
\spaceskip=5pt \xspaceskip=7pt
\begin{tikzcd} \text{a b. c} \arrow[r] & \the\spaceskip \end{tikzcd}
%%%% case everycr
\everycr={\noalign{\hrule}}
\begin{tikzcd} A \arrow[r] & B \end{tikzcd}
%%%% case everycr-array
\everycr={\noalign{\hrule}}
\begin{tikzcd} \begin{array}{c} x \\ y \end{array} \arrow[r] & B \end{tikzcd}
%%%% case everypar-arrow
\everypar={\arrow[r]}
\noindent
\begin{tikzcd} \parbox{1cm}{x} & B \end{tikzcd}
%%%% case everypar-text
\everypar={Q}
\noindent
\begin{tikzcd} \parbox{1cm}{x} \arrow[r] & B \end{tikzcd}
%%%% case everyvbox-arrow
\everyvbox={\arrow[r]}
\begin{tikzcd} \vbox{\hbox{x}} & B \end{tikzcd}
%%%% case everyvbox-text
\everyvbox={\hbox{Q}}
\begin{tikzcd} \vbox{\hbox{x}} \arrow[r] & B \end{tikzcd}
%%%% case everyhbox
\everyhbox={Q}
\begin{tikzcd} A \arrow[r] & B \end{tikzcd}
%%%% case everymath-text
\everymath={Q}
\begin{tikzcd} A \arrow[r] & B \end{tikzcd}
%%%% case everymath-style
\everymath={\displaystyle}
\begin{tikzcd} \sum_x A \arrow[r, "\sum_y"] & B \end{tikzcd}
%%%% case everymath-macro
\newcommand\h{\displaystyle}
\everymath=\expandafter{\h}
\begin{tikzcd} \sum_x A \arrow[r] & B \end{tikzcd}
%%%% case everymath-arrow
\everymath={\arrow[r]}
\begin{tikzcd} A & B \end{tikzcd}
%%%% case everydisplay
\everydisplay={\arrow[r]}
\begin{tikzcd} A \arrow[r] & B \end{tikzcd}
%%%% case everydisplay-cell
\everydisplay={Q}
\begin{tikzcd} A \arrow[r] & \text{$$x$$} \end{tikzcd}
%%%% case everyjob-eof
\everyeof={\arrow[r]}
\begin{tikzcd} A & B \end{tikzcd}
%%%% case frozen-relax
\def\h{\relax}
\begin{tikzcd} A \h \arrow[r] & B \end{tikzcd}
%%%% case math-group-limit
\DeclareSymbolFont{extra}{OT1}{cmr}{m}{n}
\DeclareSymbolFont{extrb}{OT1}{cmss}{m}{n}
\DeclareSymbolFont{extrc}{OT1}{cmtt}{m}{n}
\DeclareMathAlphabet\ma{OT1}{cmss}{bx}{n}
\DeclareMathAlphabet\mb{OT1}{cmtt}{m}{it}
\DeclareMathAlphabet\mc{OT1}{cmr}{bx}{it}
\begin{tikzcd} \ma{A} \arrow[r] & \mb{B} \mc{C} \end{tikzcd}
%%%% case mathsurround
\mathsurround=5pt
\begin{tikzcd} A \arrow[r] & B \end{tikzcd}
%%%% case mathsurround-text
\begin{tikzcd} \text{\mathsurround=5pt $x$} \arrow[r] & B \end{tikzcd}
%%%% case thinmuskip
\thinmuskip=20mu
\begin{tikzcd} A \, B \mathop{\rm x} C \arrow[r] & B \end{tikzcd}
%%%% case delimiterfactor
\delimiterfactor=500 \delimitershortfall=30pt
\begin{tikzcd} \left( \frac{a}{b} \right) \arrow[r] & B \end{tikzcd}
%%%% case nulldelimiterspace
\nulldelimiterspace=10pt
\begin{tikzcd} \frac{a}{b} \arrow[r] & B \end{tikzcd}
%%%% case scriptspace
\scriptspace=10pt
\begin{tikzcd} A^b \arrow[r] & B \end{tikzcd}
%%%% case binoppenalty
\binoppenalty=5 \relpenalty=7
\begin{tikzcd} A + B = C \arrow[r] & B \end{tikzcd}
%%%% case fam-default
\fam=1
\begin{tikzcd} 1 + B \arrow[r] & \mathrm{x} \end{tikzcd}
%%%% case skewchar
\skewchar\textfont1=`A
\begin{tikzcd} \hat{A} \arrow[r] & B \end{tikzcd}
%%%% case hyphenchar
\hyphenchar\font=`A
\begin{tikzcd} \text{a-b} \arrow[r] & B \end{tikzcd}
%%%% case uchyph
\uchyph=0 \exhyphenpenalty=3 \hyphenpenalty=4 \lefthyphenmin=1
\begin{tikzcd} \parbox{1cm}{Antidisestablishmentarianism} \arrow[r] & B \end{tikzcd}
%%%% case pretolerance
\pretolerance=-1 \tolerance=1 \emergencystretch=10pt
\begin{tikzcd} \parbox{1cm}{a b c d e f g} \arrow[r] & B \end{tikzcd}
%%%% case sfcode
\frenchspacing
\begin{tikzcd} \text{a. b} \arrow[r] & B \end{tikzcd}
%%%% case lastlinefit
\lastlinefit=500
\begin{tikzcd} \parbox{2cm}{a b c d e f g h i} \arrow[r] & B \end{tikzcd}
%%%% case pdfadjustspacing
\pdfadjustspacing=2 \pdfprotrudechars=2
\begin{tikzcd} \text{a b} \arrow[r] & B \end{tikzcd}
%%%% case predisplaysize
\begin{tikzcd} \parbox{3cm}{a\[x\]b} \arrow[r] & B \end{tikzcd}
